\section{Costly Proof and Observable Awards}
\label{sec:setting_layers}

\subsection{Screening, Investigation, and Proof}
\label{sec:setting_screening}

Antitrust agencies rarely begin with the evidence they need to
prove a procurement cartel. They begin with coarse records,
complaints, leniency leads, procurement anomalies, or repeated
interactions that make some firms and markets worth a closer
look. Legal proof comes later, after the agency recovers bid
histories, communications, leniency testimony, and other
case-specific evidence sufficient to support adjudication. The
first institutional problem is therefore not classification once
rich evidence is available; it is allocation of costly
proof-producing effort.

That distinction defines the role of an award-layer screen. The
relevant performance criterion is not whether the screen
establishes liability, but whether it improves the allocation of
costly forensic attention. Bid-distribution methods ask whether
the submitted bids look suspicious once per-bidder bid amounts
are available. An award-layer screen asks an earlier question:
before recovering those bid amounts, can the agency rank the
firms and procurement environments where bid-layer forensics is
most likely to pay off? Bid-distribution forensics remains
indispensable at the evidentiary stage; award-layer triage helps
decide where that stage should start.

The institutional sequencing problem has a substantial
law-and-economics tradition. The Becker--Stigler framework
treats enforcement intensity as a choice variable under costly
observation \citep{becker1968crime,stigler1970optimum}.
Subsequent work formalizes screening as a first-stage filter
that allocates investigative effort over a population of
potentially infringing firms
\citep{harrington2008detecting,baker2002case}: when proof is
expensive and conduct is partially observable, an inexpensive
ranking step on routinely available data is the rational
starting point of investigation. The award-layer construction
in this paper instantiates that filter. It scores firms using
information the agency already has, before deciding whether to
mobilize the costly bid-layer evidentiary stage.

What counts as proof in a procurement-cartel case is more
substantial. The agency must establish a coordinated
bid-rigging arrangement that produced anticompetitive harm,
typically through evidence of communication, allocation rules,
or bid patterns inconsistent with independent decision-making.
Brazilian competition law treats such conduct as an infraction
under Lei~$12.529/2011$; comparable frameworks in the United
States and the European Union also place the agreement element
and its effect at the center of the case
\citep{hovenkamp2005federal}. The screening object analyzed here
does none of this. It produces a ranking that helps the agency
decide where to seek the evidence those legal frameworks
require.

\subsection{The BEC Enforcement Environment}
\label{sec:setting_bec}

São Paulo's BEC platform is the right environment for this
design because it places the routine record layer on which
triage can run, the forensic-recoverable bid records, and
external CADE adjudication labels in the same procurement
system. It covers $\valSampleN$ tender-items during
$\valYearStart$--$\valYearEnd$ and preserves a contract-award
envelope for each tender-item. Procurement-cartel cases are
adjudicated by CADE, the Brazilian competition authority,
supplying external legal anchors for validation. The
combination lets us study the paper's central allocation
problem: whether information available before bid recovery can
rank where costly forensic effort should begin.

Two procurement modalities dominate. \emph{Convite} is a
sealed-bid invitation under Lei~$8.666/93$ with a statutory
three-bidder quorum and a low value cap; \emph{Pregão} is a
reverse electronic auction under Lei~$\valLeiPregao$ with
progressively lower real-time offers and no minimum-bidder rule.
The modalities differ in bidding rule and observability regime:
convite is closer to an award-only environment, while pregão
generates a richer real-time interaction footprint. We treat
this contrast as scope information, not causal identification.
The common feature matters more for the architecture: both
modalities leave an award layer that can be queried before
bid-layer evidence is recovered.

Cartel enforcement in Brazilian procurement is divided across
institutions. CADE, the federal antitrust authority,
investigates and adjudicates anticompetitive conduct including
bid rigging. Federal and state audit courts---TCU at the federal
level, TCE-SP at the state level---supervise procurement
regularity and refer cases to CADE or to the Public
Prosecutor's Office (\emph{Ministério Público}) when collusion
is suspected. The screen analyzed here lives in the
institutional space of CADE's bid-rigging caseload: it ranks
firms whose award-layer behavior warrants the costly bid-layer
recovery a CADE investigation would require. The same record
layer is, however, already routinely consumed by audit-court and
prosecutorial staff during procurement oversight
\citep{oecd2019government}, which is one reason the architecture
is institutionally implementable rather than purely academic.

\subsection{Award Layer and Forensic-Recoverable Bid Layer}
\label{sec:setting_layers_two}

BEC exposes two analytically distinct layers. The
\emph{award layer} carries the fields an enforcement analyst can
query routinely: participant identity, winner identity, item
code, negotiated price, procuring unit, year, and procurement
modality. The \emph{forensic-recoverable bid layer} carries
per-bidder bid amounts, bid timestamps, and bid-revision
sequences. Those bid records exist in the BEC operational system,
but they are not exported to the ordinary analytical warehouse.
They must be recovered through an administrative LANCES export:
feasible for evaluation and case work, but costly relative to
querying the award layer. OECD reviews of Brazilian procurement
infrastructure describe this kind of administrative latency as a
practical constraint on routine monitoring
\citep{oecd2019government}. The constraint is not data
non-existence; it is the cost and timing of turning operational
bid files into analyzable forensic evidence.

Figure~\ref{fig:coarsening_v18} is not an accuracy horse race; it
is an information-cost diagram. The left side reports the
bid-layer features required by the Imhof--Wallimann-style forensic
benchmark: within-tender coefficient of variation, dispersion
across tenders, skewness, kurtosis, normalized spread, log
min--max, and second-lowest bid distance. The right side reports
the four fields used by the award-layer screen. The empirical
question is whether the right-hand layer can rank where the
left-hand layer should be opened first; the bid-layer features
remain the forensic-stage benchmark, recovered through LANCES
only after triage has selected where to look.

\begin{figure}[htbp]
\centering
\includegraphics[width=0.92\textwidth]{./output/figures/fig_data_coarsening}
\caption[Information layers in procurement-cartel enforcement]{Information
layers in procurement-cartel enforcement. Bid-distribution
forensics requires per-bidder bid amounts and within-tender bid
moments; the award-layer screen uses participant identity, winner
identity, item code, and negotiated price. The empirical question
is whether the right-hand layer can rank where the left-hand layer
should be recovered first.}
\label{fig:coarsening_v18}
\end{figure}

The same institutional distinction appears outside BEC. In the
United States, FPDS-NG preserves federal contract awards as
routine administrative records, while bid files sit with
contracting agencies and are typically recovered through
case-specific process. In the European Union, TED mandates
publication of contract-award notices, but per-bidder bid
amounts are preserved heterogeneously across member-state
systems; a recent procurement-screening review reports that
fewer than half of surveyed systems make bid amounts retrievable
at the analytical-warehouse layer
\citep{sanchezgraells2019screening}. These comparisons are not
external validation of the BEC estimates. They establish a
narrower point: the institutional distinction between routine
award records and costly bid files is recognizable outside this
setting.

The bid-distribution literature presumes the left-hand side of
Figure~\ref{fig:coarsening_v18}. The variance and dispersion
screens of \citet{abrantesmetz2006a} use within-tender bid
moments; the machine-learning extensions of
\citet{imhof2018screening,imhof2019detecting,huber2019machine,wallimann2023machine}
apply richer functions to the same per-bidder bid amounts. The
coalition-inference apparatus of \citet{conley2016detecting}
requires bid-level histories. These methods are not portable to
the right-hand side because their inputs are not yet recovered
there, not because the underlying logic fails. The paper's
question is therefore not whether to use bid-distribution
moments---we use them at the forensic stage---but whether
award-layer information can decide where to recover those
moments first.

\subsection{CADE Labels and the Cartel-Adjacency Target}
\label{sec:setting_cade}

CADE adjudications provide the legal anchors for the validation
exercise. CADE enforces Brazilian competition law under
Lei~$12.529/2011$, and its procurement-cartel decisions identify
firms legally named in cartel cases. The portfolio linked to
BEC contains $12$ adjudicated procurement-cartel cases with $65$
firm-defendants, of which $\valDirectCADE$ are active in BEC
during the sample window. Final decisions span
$2015$--$2025$, with process-to-ruling intervals of $4$--$16$
years. That lag is central to the screening interpretation: an
award-layer screen has value only if it can move suspicion
toward forensic investigation before the completed adjudication
record exists.

The adjudication process behind those labels has a standard
structure. An investigation is opened by CADE's
\emph{Superintend\^encia-Geral} (SG), often following a
leniency application, a \emph{Termo de Compromisso de Cessa\c{c}\~ao}
settlement, or a referral from another enforcement body. The
SG issues a recommendation; CADE's Tribunal reviews and rules.
Cartel cases are heavily leniency-driven: the leniency
provisions of Lei~$12.529/2011$, operationalized through CADE
regulations, supply most of the direct-defendant identifications
in our portfolio. This matters for what direct defendants are.
They are firms whose membership was established through legal
process, often via the cooperation of another cartelist. The
award-layer screen has no equivalent legal authority and
produces no analogous identification.

The adjudication record creates two distinct empirical objects.
Direct CADE defendants are legal membership anchors: they
locate cartel-affected environments but are not the natural
target of a loser-side award-layer statistic, because in
procurement cartels they often occupy winner-side or allocation
roles. Cobidders are different. They are always-loser firms
that appear alongside direct defendants in BEC tender-items;
they are indirect labels, and we do not treat them as proof of
cartel membership. They are the loser-side adjacency object the
award layer can plausibly recover: firms that lose, continue to
appear, and are observed near legally adjudicated cartel
anchors.

This distinction is not a fallback after failed defendant
classification. It follows from the information layer. Award
records contain participation and outcome information; they do
not contain agreements, communications, or the tender-specific
conduct needed for liability. A statistic built from zero wins
and repeated participation should therefore be evaluated against
loser-side exposure, not legal membership. Weak performance
against direct defendants is a scope check rather than a
failure: it shows that the statistic is not a generic
cartel-membership detector.

The rest of the analysis follows that boundary. We evaluate
whether award records rank cartel-adjacent loser-side firms well
enough to allocate bid-layer forensic effort. The screen remains
a triage device; it does not adjudicate liability.
